% Options for packages loaded elsewhere
% Options for packages loaded elsewhere
\PassOptionsToPackage{unicode}{hyperref}
\PassOptionsToPackage{hyphens}{url}
\PassOptionsToPackage{dvipsnames,svgnames,x11names}{xcolor}
%
\documentclass[
  letterpaper,
  DIV=11,
  numbers=noendperiod]{scrreprt}
\usepackage{xcolor}
\usepackage{amsmath,amssymb}
\setcounter{secnumdepth}{5}
\usepackage{iftex}
\ifPDFTeX
  \usepackage[T1]{fontenc}
  \usepackage[utf8]{inputenc}
  \usepackage{textcomp} % provide euro and other symbols
\else % if luatex or xetex
  \usepackage{unicode-math} % this also loads fontspec
  \defaultfontfeatures{Scale=MatchLowercase}
  \defaultfontfeatures[\rmfamily]{Ligatures=TeX,Scale=1}
\fi
\usepackage{lmodern}
\ifPDFTeX\else
  % xetex/luatex font selection
\fi
% Use upquote if available, for straight quotes in verbatim environments
\IfFileExists{upquote.sty}{\usepackage{upquote}}{}
\IfFileExists{microtype.sty}{% use microtype if available
  \usepackage[]{microtype}
  \UseMicrotypeSet[protrusion]{basicmath} % disable protrusion for tt fonts
}{}
\makeatletter
\@ifundefined{KOMAClassName}{% if non-KOMA class
  \IfFileExists{parskip.sty}{%
    \usepackage{parskip}
  }{% else
    \setlength{\parindent}{0pt}
    \setlength{\parskip}{6pt plus 2pt minus 1pt}}
}{% if KOMA class
  \KOMAoptions{parskip=half}}
\makeatother
% Make \paragraph and \subparagraph free-standing
\makeatletter
\ifx\paragraph\undefined\else
  \let\oldparagraph\paragraph
  \renewcommand{\paragraph}{
    \@ifstar
      \xxxParagraphStar
      \xxxParagraphNoStar
  }
  \newcommand{\xxxParagraphStar}[1]{\oldparagraph*{#1}\mbox{}}
  \newcommand{\xxxParagraphNoStar}[1]{\oldparagraph{#1}\mbox{}}
\fi
\ifx\subparagraph\undefined\else
  \let\oldsubparagraph\subparagraph
  \renewcommand{\subparagraph}{
    \@ifstar
      \xxxSubParagraphStar
      \xxxSubParagraphNoStar
  }
  \newcommand{\xxxSubParagraphStar}[1]{\oldsubparagraph*{#1}\mbox{}}
  \newcommand{\xxxSubParagraphNoStar}[1]{\oldsubparagraph{#1}\mbox{}}
\fi
\makeatother

\usepackage{color}
\usepackage{fancyvrb}
\newcommand{\VerbBar}{|}
\newcommand{\VERB}{\Verb[commandchars=\\\{\}]}
\DefineVerbatimEnvironment{Highlighting}{Verbatim}{commandchars=\\\{\}}
% Add ',fontsize=\small' for more characters per line
\usepackage{framed}
\definecolor{shadecolor}{RGB}{241,243,245}
\newenvironment{Shaded}{\begin{snugshade}}{\end{snugshade}}
\newcommand{\AlertTok}[1]{\textcolor[rgb]{0.68,0.00,0.00}{#1}}
\newcommand{\AnnotationTok}[1]{\textcolor[rgb]{0.37,0.37,0.37}{#1}}
\newcommand{\AttributeTok}[1]{\textcolor[rgb]{0.40,0.45,0.13}{#1}}
\newcommand{\BaseNTok}[1]{\textcolor[rgb]{0.68,0.00,0.00}{#1}}
\newcommand{\BuiltInTok}[1]{\textcolor[rgb]{0.00,0.23,0.31}{#1}}
\newcommand{\CharTok}[1]{\textcolor[rgb]{0.13,0.47,0.30}{#1}}
\newcommand{\CommentTok}[1]{\textcolor[rgb]{0.37,0.37,0.37}{#1}}
\newcommand{\CommentVarTok}[1]{\textcolor[rgb]{0.37,0.37,0.37}{\textit{#1}}}
\newcommand{\ConstantTok}[1]{\textcolor[rgb]{0.56,0.35,0.01}{#1}}
\newcommand{\ControlFlowTok}[1]{\textcolor[rgb]{0.00,0.23,0.31}{\textbf{#1}}}
\newcommand{\DataTypeTok}[1]{\textcolor[rgb]{0.68,0.00,0.00}{#1}}
\newcommand{\DecValTok}[1]{\textcolor[rgb]{0.68,0.00,0.00}{#1}}
\newcommand{\DocumentationTok}[1]{\textcolor[rgb]{0.37,0.37,0.37}{\textit{#1}}}
\newcommand{\ErrorTok}[1]{\textcolor[rgb]{0.68,0.00,0.00}{#1}}
\newcommand{\ExtensionTok}[1]{\textcolor[rgb]{0.00,0.23,0.31}{#1}}
\newcommand{\FloatTok}[1]{\textcolor[rgb]{0.68,0.00,0.00}{#1}}
\newcommand{\FunctionTok}[1]{\textcolor[rgb]{0.28,0.35,0.67}{#1}}
\newcommand{\ImportTok}[1]{\textcolor[rgb]{0.00,0.46,0.62}{#1}}
\newcommand{\InformationTok}[1]{\textcolor[rgb]{0.37,0.37,0.37}{#1}}
\newcommand{\KeywordTok}[1]{\textcolor[rgb]{0.00,0.23,0.31}{\textbf{#1}}}
\newcommand{\NormalTok}[1]{\textcolor[rgb]{0.00,0.23,0.31}{#1}}
\newcommand{\OperatorTok}[1]{\textcolor[rgb]{0.37,0.37,0.37}{#1}}
\newcommand{\OtherTok}[1]{\textcolor[rgb]{0.00,0.23,0.31}{#1}}
\newcommand{\PreprocessorTok}[1]{\textcolor[rgb]{0.68,0.00,0.00}{#1}}
\newcommand{\RegionMarkerTok}[1]{\textcolor[rgb]{0.00,0.23,0.31}{#1}}
\newcommand{\SpecialCharTok}[1]{\textcolor[rgb]{0.37,0.37,0.37}{#1}}
\newcommand{\SpecialStringTok}[1]{\textcolor[rgb]{0.13,0.47,0.30}{#1}}
\newcommand{\StringTok}[1]{\textcolor[rgb]{0.13,0.47,0.30}{#1}}
\newcommand{\VariableTok}[1]{\textcolor[rgb]{0.07,0.07,0.07}{#1}}
\newcommand{\VerbatimStringTok}[1]{\textcolor[rgb]{0.13,0.47,0.30}{#1}}
\newcommand{\WarningTok}[1]{\textcolor[rgb]{0.37,0.37,0.37}{\textit{#1}}}

\usepackage{longtable,booktabs,array}
\newcounter{none} % for unnumbered tables
\usepackage{calc} % for calculating minipage widths
% Correct order of tables after \paragraph or \subparagraph
\usepackage{etoolbox}
\makeatletter
\patchcmd\longtable{\par}{\if@noskipsec\mbox{}\fi\par}{}{}
\makeatother
% Allow footnotes in longtable head/foot
\IfFileExists{footnotehyper.sty}{\usepackage{footnotehyper}}{\usepackage{footnote}}
\makesavenoteenv{longtable}
\usepackage{graphicx}
\makeatletter
\newsavebox\pandoc@box
\newcommand*\pandocbounded[1]{% scales image to fit in text height/width
  \sbox\pandoc@box{#1}%
  \Gscale@div\@tempa{\textheight}{\dimexpr\ht\pandoc@box+\dp\pandoc@box\relax}%
  \Gscale@div\@tempb{\linewidth}{\wd\pandoc@box}%
  \ifdim\@tempb\p@<\@tempa\p@\let\@tempa\@tempb\fi% select the smaller of both
  \ifdim\@tempa\p@<\p@\scalebox{\@tempa}{\usebox\pandoc@box}%
  \else\usebox{\pandoc@box}%
  \fi%
}
% Set default figure placement to htbp
\def\fps@figure{htbp}
\makeatother





\setlength{\emergencystretch}{3em} % prevent overfull lines

\providecommand{\tightlist}{%
  \setlength{\itemsep}{0pt}\setlength{\parskip}{0pt}}



 


\KOMAoption{captions}{tableheading}
\makeatletter
\@ifpackageloaded{bookmark}{}{\usepackage{bookmark}}
\makeatother
\makeatletter
\@ifpackageloaded{caption}{}{\usepackage{caption}}
\AtBeginDocument{%
\ifdefined\contentsname
  \renewcommand*\contentsname{Table of contents}
\else
  \newcommand\contentsname{Table of contents}
\fi
\ifdefined\listfigurename
  \renewcommand*\listfigurename{List of Figures}
\else
  \newcommand\listfigurename{List of Figures}
\fi
\ifdefined\listtablename
  \renewcommand*\listtablename{List of Tables}
\else
  \newcommand\listtablename{List of Tables}
\fi
\ifdefined\figurename
  \renewcommand*\figurename{Figure}
\else
  \newcommand\figurename{Figure}
\fi
\ifdefined\tablename
  \renewcommand*\tablename{Table}
\else
  \newcommand\tablename{Table}
\fi
}
\@ifpackageloaded{float}{}{\usepackage{float}}
\floatstyle{ruled}
\@ifundefined{c@chapter}{\newfloat{codelisting}{h}{lop}}{\newfloat{codelisting}{h}{lop}[chapter]}
\floatname{codelisting}{Listing}
\newcommand*\listoflistings{\listof{codelisting}{List of Listings}}
\makeatother
\makeatletter
\makeatother
\makeatletter
\@ifpackageloaded{caption}{}{\usepackage{caption}}
\@ifpackageloaded{subcaption}{}{\usepackage{subcaption}}
\makeatother
\usepackage{bookmark}
\IfFileExists{xurl.sty}{\usepackage{xurl}}{} % add URL line breaks if available
\urlstyle{same}
\hypersetup{
  pdftitle={21CSC106T-Foundation of Data Sience},
  pdfauthor={Jubilant J},
  colorlinks=true,
  linkcolor={blue},
  filecolor={Maroon},
  citecolor={Blue},
  urlcolor={Blue},
  pdfcreator={LaTeX via pandoc}}


\title{21CSC106T-Foundation of Data Sience}
\author{Jubilant J}
\date{2026-01-05}
\begin{document}
\maketitle

\renewcommand*\contentsname{Table of contents}
{
\setcounter{tocdepth}{2}
\tableofcontents
}

\bookmarksetup{startatroot}

\chapter*{Preface}\label{preface}
\addcontentsline{toc}{chapter}{Preface}

\markboth{Preface}{Preface}

This is a Quarto book.

To learn more about Quarto books visit
\url{https://quarto.org/docs/books}.

\bookmarksetup{startatroot}

\chapter{Introduction}\label{introduction}

\subsection{Topics Covered}\label{topics-covered}

An Introduction to Data Analysis - Data Analysis - Knowledge domains of
Data Analyst: Computer Science, Mathematics, and statistics - Machine
Learning \& AI, Professional fields of Application - Introduction to
Data - Understanding the nature of Data - Data -- Information;
Information -- Knowledge - Types of Data - Data Analysis Process -
Quantitative Data Analysis - Qualitative Data Analysis - Python -- The
Programming Language - Python 2 and Python 3 - Python Package Index -
IDEs for python - Scipy: Numpy- Pandas, Matplotlib.

\section{An Introduction to Data
Analysis}\label{an-introduction-to-data-analysis}

Data analysis starts with identifying a problem that can be solved with
data. Once the problem is identified, data is collected, cleaned,
processed, and analyzed. The purpose of analyzing the data is to
identify trends, patterns, and meaningful insights, with the ultimate
goal of solving the original problem.

The following five-step framework helps to analyze the data :

\begin{enumerate}
\def\labelenumi{\arabic{enumi}.}
\tightlist
\item
  Identify business questions
\item
  Collect and store data
\item
  Clean and prepare data
\item
  Analyze data
\item
  Visualize and communicate data.
\end{enumerate}

\subsection{Data Analysis Process}\label{data-analysis-process}

Data analysis is schematized as a process chain consisting of the
following sequence of stages:

\begin{itemize}
\tightlist
\item
  \textbf{Problem definition}: Data analysis always starts with a
  problem to be solved. Requires focusing the system, planning, and
  choosing an effective, cross-disciplinary team.
\item
  \textbf{Data extraction}: Obtaining the data to perform the analysis.
  Data selection is crucial.
\item
  \textbf{Data preparation}: Cleaning, normalizing, and transforming
  data into an optimized dataset. Handling invalid, ambiguous, or
  missing values {[}7{]}.
\item
  \textbf{Data exploration and visualization}: Searching data
  graphically or statistically to find patterns, connections, and
  relationships. Includes summarizing, grouping, and identifying trends
  {[}7, 8{]}.
\item
  \textbf{Predictive modelling}: Creating a statistical model to predict
  results (Regression models for numeric data, Classification/Clustering
  models for descriptive/categorical data) {[}8, 9{]}.
\item
  \textbf{Model validation/test}: Validating the model built on the
  basis of starting data (e.g., using Cross Validation techniques) {[}9,
  10{]}.
\item
  \textbf{Deployment}: The final step which aims to present the results.
  Involves documenting analysis results, decision deployment, risk
  analysis, and measuring the business impact {[}10, 11{]}.
\end{itemize}

\section{Knowledge Domains of Data
Analyst}\label{knowledge-domains-of-data-analyst}

Data Science comprises three main subject areas {[}12, 13{]}:

\begin{enumerate}
\def\labelenumi{\arabic{enumi}.}
\tightlist
\item
  \textbf{Computer Science and Programming}: The study of computational
  tools like programming languages and software libraries.
\item
  \textbf{Statistics and Machine Learning}: Forms the theoretical
  foundations of data science methods and algorithms.
\item
  \textbf{Domain Knowledge}: The general discipline or field to which
  data science is applied (e.g., medicine, finance, biotech). It is
  required to interpret results properly {[}12, 13{]}.
\end{enumerate}

\section{Machine Learning \& AI, Professional Fields of
Application}\label{machine-learning-ai-professional-fields-of-application}

\begin{itemize}
\tightlist
\item
  \textbf{Image Recognition}: Cataloging and detecting objects, face
  detection, and pattern recognition {[}14, 15{]}.
\item
  \textbf{Speech Recognition}: Representing speech signals with numbers
  (e.g., Amazon's Alexa, Apple's Siri) {[}15{]}.
\item
  \textbf{Predict Traffic Patterns}: Analyzing present traffic data to
  generate predictions and identify the fastest route {[}15, 16{]}.
\item
  \textbf{E-commerce Product Recommendations}: Tracking customer
  behavior to recommend products {[}16{]}.
\item
  \textbf{Self-Driving Cars}: Utilizing unsupervised learning algorithms
  via cameras and sensors {[}17{]}.
\item
  \textbf{Catching Email Spam}: Filtering incoming emails using ML
  algorithms {[}17, 18{]}.
\item
  \textbf{Catching Malware}: Analyzing suspicious activities in
  environments like Android {[}18{]}.
\item
  \textbf{Virtual Personal Assistant}: Accessing relevant information
  via text or voice using Natural Language Processing {[}19{]}.
\item
  \textbf{Online Fraud Recognition}: Examining user profiles for unusual
  patterns {[}19, 20{]}.
\item
  \textbf{Stock Market and Day Trading}: Automating investment
  activities using algorithmic trading {[}20{]}.
\end{itemize}

\section{Quantitative and Qualitative Data
Analysis}\label{quantitative-and-qualitative-data-analysis}

\begin{itemize}
\tightlist
\item
  \textbf{Quantitative}: Analysed data (Numerical or Categorical),
  allows mathematical processing, and provides more objective
  conclusions {[}21{]}.
\item
  \textbf{Qualitative}: Values expressed through descriptions in Natural
  Language (text, visual, or audio data), leading to subjective
  interpretations {[}21{]}.
\end{itemize}

\subsection{Open Data}\label{open-data}

Freely available data sources on the internet supporting the growing
demand for data (e.g., DataHub, WHO, Data.gov, EU Open Data Portal, AWS
public datasets, Google Trends) {[}21, 22{]}.

\section{Understanding the Nature of
Data}\label{understanding-the-nature-of-data}

\begin{itemize}
\tightlist
\item
  \textbf{Data (Raw)}: Raw alphanumeric values obtained through
  different acquisition methods {[}23{]}.
\item
  \textbf{Information (Processed)}: Created when data is processed,
  organized, or structured to provide context and meaning {[}23{]}.
\item
  \textbf{Knowledge (Actionable)}: Unique to each individual; it is the
  accumulation of past experience and insight that shapes how we
  interpret information to make decisions {[}23{]}.
\end{itemize}

\subsection{Types of Data}\label{types-of-data}

\begin{enumerate}
\def\labelenumi{\arabic{enumi}.}
\tightlist
\item
  \textbf{Categorical or Qualitative Data} {[}24{]}

  \begin{itemize}
  \tightlist
  \item
    \textbf{Nominal Data}: Represents discrete units, cannot be ordered
    or measured (e.g., Gender, Nationality, Hair color). Uses one-hot
    encoding {[}24, 25{]}.
  \item
    \textbf{Ordinal Data}: Indicates an ordered or graded relationship
    (e.g., Ranks, Opinion scales, Socioeconomic status) {[}25, 26{]}.
  \end{itemize}
\item
  \textbf{Numerical or Quantitative Data} {[}24{]}

  \begin{itemize}
  \tightlist
  \item
    \textbf{Discrete Data}: Can take only finite discrete values,
    counted in whole numbers (e.g., Number of students) {[}26, 27{]}.
  \item
    \textbf{Continuous Data}: Data that can be calculated with an
    infinite number of probable values (e.g., Temperature range)
    {[}27{]}.
  \end{itemize}
\end{enumerate}

\section{Python and Data Analysis}\label{python-and-data-analysis}

Python is widely used in scientific circles because of its large number
of libraries for analysis and data manipulation {[}27{]}. *
\textbf{Python Package Index (PyPI)}: A repository of software for the
Python programming language helping you find and install packages (using
\texttt{pip}) {[}28, 29{]}. * \textbf{Characteristics}: Interpreted,
Portable, Object-oriented, Interactive, Interfaced, Open source
{[}30{]}.

\subsection{Interpreters vs Compilers}\label{interpreters-vs-compilers}

\begin{itemize}
\tightlist
\item
  \textbf{Compiler}: Takes entire program as input, generates
  intermediate object code, executes faster, requires more memory.
  (e.g., C compiler) {[}31, 32{]}.
\item
  \textbf{Interpreter}: Takes single instruction as input, generates no
  object code, is slower in execution, and uses less memory. Errors are
  displayed line by line. (e.g., Python) {[}31, 32{]}.
\end{itemize}

\subsection{Python Implementations}\label{python-implementations}

\begin{itemize}
\tightlist
\item
  \textbf{Cython}: Translates Python code into C {[}32{]}.
\item
  \textbf{Jython}: An implementation of Python built and compiled in
  Java {[}33{]}.
\item
  \textbf{IronPython}: Developed in C\# (for Windows) {[}32{]}.
\item
  \textbf{PyPy}: A JIT (just-in-time) compiler that converts Python
  directly to machine code at runtime {[}33{]}.
\end{itemize}

\section{Python Basics}\label{python-basics}

\begin{itemize}
\tightlist
\item
  \textbf{Comments}: Used to mark non-executable code. Single-line
  (\texttt{\#}) and Multi-line Docstrings
  (\texttt{\textquotesingle{}\textquotesingle{}\textquotesingle{}} or
  \texttt{"""}) {[}34, 35{]}.
\item
  \textbf{Identifiers}: Names used for variables, functions, or classes.
  Case-sensitive. Cannot start with a digit or be a keyword. Allowed
  special symbol is underscore \texttt{\_} {[}36-38{]}.
\item
  \textbf{Keywords}: Reserved words with special meaning (e.g.,
  \texttt{False}, \texttt{class}, \texttt{finally}, \texttt{is},
  \texttt{return}). Python 3.5 has 33 keywords {[}39{]}.
\item
  \textbf{Variables}: Reserved memory locations. No explicit declaration
  needed; created upon assignment using the \texttt{=} operator {[}40,
  41{]}.

  \begin{itemize}
  \tightlist
  \item
    \emph{Chained Assignment}: \texttt{a\ =\ b\ =\ c\ =\ 10} {[}42{]}
  \item
    \emph{Multiple Assignment}: \texttt{value1,\ value2\ =\ 1,\ 2.5}
    {[}43{]}
  \end{itemize}
\item
  \textbf{Expressions}: A combination of values, variables, and
  operators {[}44{]}.
\item
  \textbf{Augmented Assignment}: Shorthand operators like \texttt{+=},
  \texttt{-=}, \texttt{*=}, \texttt{/=}, \texttt{\%=} (e.g.,
  \texttt{a\ +=\ b} is equivalent to \texttt{a\ =\ a\ +\ b}) {[}45,
  46{]}.
\end{itemize}

\section{Data Structures in Python}\label{data-structures-in-python}

\subsection{Dictionary}\label{dictionary}

Unordered collection of key-value pairs. Keys must be unique and
immutable {[}47-49{]}.

\begin{Shaded}
\begin{Highlighting}[]
\BuiltInTok{dict} \OperatorTok{=}\NormalTok{ \{}\StringTok{\textquotesingle{}name\textquotesingle{}}\NormalTok{:}\StringTok{\textquotesingle{}William\textquotesingle{}}\NormalTok{, }\StringTok{\textquotesingle{}age\textquotesingle{}}\NormalTok{:}\DecValTok{25}\NormalTok{, }\StringTok{\textquotesingle{}city\textquotesingle{}}\NormalTok{:}\StringTok{\textquotesingle{}London\textquotesingle{}}\NormalTok{\}}
\BuiltInTok{print}\NormalTok{(}\BuiltInTok{dict}\NormalTok{[}\StringTok{"name"}\NormalTok{]) }\CommentTok{\# Outputs: \textquotesingle{}William\textquotesingle{}}
\end{Highlighting}
\end{Shaded}

\begin{verbatim}
William
\end{verbatim}

\bookmarksetup{startatroot}

\chapter{Module 2}\label{module-2}

\bookmarksetup{startatroot}

\chapter{Module-3}\label{module-3}

\bookmarksetup{startatroot}

\chapter{Module-4}\label{module-4}

\bookmarksetup{startatroot}

\chapter{Module-5}\label{module-5}

\bookmarksetup{startatroot}

\chapter{Appendix}\label{appendix}

\bookmarksetup{startatroot}

\chapter{Lab}\label{lab}

\bookmarksetup{startatroot}

\chapter{Summary}\label{summary}

In summary, this book has no content whatsoever.

\bookmarksetup{startatroot}

\chapter*{References}\label{references}
\addcontentsline{toc}{chapter}{References}

\markboth{References}{References}

\protect\phantomsection\label{refs}




\end{document}
